\chapter{Contracts Under Uncertainty}
\label{ch:contract}

\section{Contracts as bets and as messages}

A contract allocates risk. A fixed price construction contract puts the risk of cost overruns on the builder; a cost plus contract puts it on the owner. An insurance policy transfers the risk of loss from the insured to the insurer for a premium. A futures contract fixes a price today for delivery tomorrow, so that each party bears the risk that the market moves against it. In each case the parties have beliefs about the uncertain future, and the terms reflect those beliefs and their attitudes to risk. Contracts are, among other things, bets.\index{contract law}

A contract is also a message. Its words are chosen by the parties to communicate their agreement, and when a dispute arises the court must infer what they meant. That inference is a problem of interpretation, and interpretation, we will argue, is Bayesian inference about intentions from text and context.

This chapter applies the tools of the book to four questions in English contract law: interpretation, implied terms, remoteness of damage and penalty clauses. It then turns to frustration, to the American Uniform Commercial Code and to machines. Each of the four has been the subject of leading decisions in the last thirty years, and each can be read as a judgment about probabilities and expected values.

\section{Interpretation as inference}

\subsection{The modern approach}

In \emph{Investors Compensation Scheme v West Bromwich Building Society} Lord Hoffmann restated the principles of contractual interpretation \citep{ics1998}. Interpretation is the ascertainment of the meaning which the document would convey to a reasonable person having all the background knowledge which would reasonably have been available to the parties in the situation in which they were at the time of the contract. The background, the matrix of fact, includes absolutely anything which would have affected the way in which the language would have been understood by a reasonable man, subject to exclusions such as the parties' previous negotiations. The meaning of a document is not the same as the meaning of its words: if something has gone wrong with the language, the law does not require judges to attribute to the parties an intention which they plainly could not have had.\index{interpretation}\index{Investors Compensation Scheme}

Later decisions qualified the emphasis on context. In \emph{Arnold v Britton} the Supreme Court held that the natural meaning of the words is of central importance, that commercial common sense should not be invoked to undervalue the language, and that the court should not rewrite a bargain that turned out badly for one party \citep{arnold2015}. In \emph{Wood v Capita Insurance Services} it described interpretation as a unitary exercise in which textualism and contextualism are tools, the balance between them depending on the circumstances, including the quality of drafting \citep{wood2017}.

\subsection{The Bayesian reading}

Let $m$ denote a possible meaning of a clause, an allocation of rights and obligations that the parties might have intended. Let $T$ be the text and $C$ the context. The court seeks the most probable meaning, or more precisely the meaning a reasonable person would attribute, which we read as the posterior mode (the single most probable meaning) or a meaning with high posterior probability:
\begin{equation}
\Prob(m\given T, C) \propto \Prob(T\given m, C)\,\Prob(m\given C).
\end{equation}
The likelihood $\Prob(T\given m, C)$ is the probability that parties intending $m$ would have written this text. A careful drafter intending $m$ would use words that express $m$ clearly, so the likelihood is high for meanings close to the natural meaning of the words and low for others. The prior $\Prob(m\given C)$ is the probability that parties in this commercial context would intend $m$, which reflects commercial common sense: meanings that make business sense have higher prior probability than meanings that would be absurd.\index{interpretation!as Bayesian inference}

This reading reconciles the cases. When the text is clear and was professionally drafted, the likelihood is sharply peaked, and it dominates the prior; the natural meaning prevails even if it seems commercially odd, as in \emph{Arnold v Britton}, where service charge clauses in leases produced escalating charges that became absurd over time but were clearly expressed. When the text is poorly drafted or ambiguous, the likelihood is flat across several meanings, and the prior, commercial common sense, decides. When something has plainly gone wrong with the language, as in Lord Hoffmann's examples, the text is not as improbable under the sensible meaning as it appears, because careful drafters make mistakes, so the likelihood ratio between the literal and the sensible meaning is modest, and the low prior of the absurd literal meaning decides. \emph{Wood v Capita}'s unitary exercise, in which the weight given to text and context depends on the quality of drafting, matches the Bayesian prescription: the precision of the likelihood depends on how carefully the text was written.

\begin{casebox}{Arnold v Britton [2015] UKSC 36}
Leases of chalets in a holiday park required lessees to pay a service charge of a fixed sum, increasing by ten per cent compound each year. When the leases were granted, at a time of high inflation, this seemed reasonable. By the 2010s, with low inflation, the charges had become large relative to the cost of services, and by the end of the leases they would be enormous. The lessees argued that the clause should be read as requiring a proportionate share of the costs, with the stated sum as a cap. The Supreme Court held that the words were clear and must be given effect, even though the result was unfavourable to the lessees. Lord Neuberger said that the purpose of interpretation is to identify what the parties agreed, not what the court thinks they should have agreed. In Bayesian terms, the text was precise, and its likelihood ratio between the literal meaning and the proposed alternative was overwhelming; the prior from commercial common sense could not overturn it, especially since the bargain had made sense when made.
\end{casebox}

\section{Implied terms}

Courts sometimes imply terms into contracts: terms the parties did not express but which the law treats as part of the agreement. In \emph{Marks and Spencer v BNP Paribas} the Supreme Court held that a term will be implied in fact only if it is necessary to give business efficacy to the contract, or so obvious that it goes without saying, and that the test is not one of reasonableness \citep{marks2015}. The officious bystander test asks whether, had a bystander suggested the term to the parties at the time of contracting, they would both have testily replied, of course.\index{implied terms}

The Bayesian reading is that a term is implied when the posterior probability that reasonable parties in their position would both have agreed to it, given the contract and its context, is close to one. The threshold is high because implying a term rewrites the bargain, and a court that implied terms whenever they were probable would substitute its judgment for the parties'. The high threshold protects the parties' freedom to leave matters unaddressed, perhaps deliberately. Reasonableness alone is not enough, because many reasonable terms are ones the parties might well have rejected.

\section{Remoteness and foreseeability}

\subsection{Hadley v Baxendale}

In \emph{Hadley v Baxendale} a mill's crankshaft broke, and the millers sent it to be copied by engineers in Greenwich, using a carrier who delayed delivery \citep{hadley1854}. The mill was idle for longer than it would have been, and the millers sued for lost profits. The court held that damages for breach of contract are recoverable only if they arise naturally from the breach, according to the usual course of things, or may reasonably be supposed to have been in the contemplation of both parties at the time they made the contract as the probable result of the breach. On the facts as the court took them, the carrier had not been told that the mill could not run without the shaft, and the loss of profits was not a natural consequence, since the millers might have had a spare shaft.\index{Hadley v Baxendale}\index{remoteness}

\subsection{Probability and the scope of assumed responsibility}

The rule in \emph{Hadley} is a probability test: a loss is recoverable if it was a probable result of breach, as contemplated by the parties. Later cases debated how probable. In \emph{The Heron II} the House of Lords required that the loss be of a kind that was not unlikely, a serious possibility, a real danger, a higher threshold than the foreseeability required in tort \citep{heron1969}. In \emph{The Achilleas} two of the Law Lords added that the question is also whether the defendant assumed responsibility for the kind of loss in question, as judged by the understanding in the market \citep{achilleas2008}. A charterer who returned a ship nine days late was liable only for the loss of the difference in hire for the nine days, not for the reduction in hire over the whole of a lucrative follow on fixture that the owners had to accept, because of the delay, to keep it. The reason was that the market understood that charterers did not assume responsibility for such losses.

The Bayesian and decision theoretic reading is this. At the time of contracting, each party has beliefs about the losses that breach might cause. The price is set on the basis of the losses the carrier expects to bear. If the carrier is liable for losses it could not foresee, it bears risks it has not priced, and in equilibrium it will charge all customers more, including those whose losses from delay are small. The rule in \emph{Hadley} gives the party with private information about unusual losses an incentive to disclose it, so that the carrier can price accordingly or take special care. Ayres and Gertner, whose term for such defaults is penalty defaults, described this as an information forcing effect \citep{ayres1989}. \emph{The Achilleas} adds that the allocation should reflect market understanding, which is the prior that a reasonable party in the market would hold about who bears which risks.

\section{Penalty clauses}

Contracts often stipulate the damages payable on breach. At common law a stipulated sum that was a genuine pre estimate of loss was enforceable as liquidated damages, while one that was a penalty, designed to deter breach rather than compensate, was not. In \emph{Cavendish Square Holding v Makdessi} the Supreme Court reformulated the test: a clause is penal if it is a secondary obligation imposing a detriment out of all proportion to any legitimate interest of the innocent party in the enforcement of the primary obligation \citep{cavendish2015}. The genuine pre estimate test was too narrow, because parties may have legitimate interests in performance beyond compensation for measurable loss.\index{penalty clauses}\index{Cavendish v Makdessi}

From an expected utility standpoint, a stipulated damages clause does several things. It reduces the cost of proving loss, which may be large and uncertain. It can signal the promisor's confidence: a builder who accepts large damages for delay signals that it expects to finish on time. And it allocates risk: a large stipulated sum transfers risk from the promisee to the promisor. A rational promisee would pay for these benefits through the price. The traditional suspicion of penalties assumed that large stipulated sums were extracted by superior bargaining power or reflected the promisor's overconfidence. \emph{Makdessi} narrowed the doctrine, recognising that commercial parties of comparable bargaining power can reasonably agree to strong incentives for performance. The remaining protection against oppressive clauses rests on proportionality to legitimate interests, which a decision theorist would assess by asking whether the clause plausibly increased the expected joint surplus of the parties at the time of contracting.

\section{Frustration and the allocation of unforeseen risk}

When an event occurs after contracting that makes performance impossible or radically different from what was contemplated, the contract may be frustrated, discharging both parties. In \emph{Taylor v Caldwell} a music hall burned down before a series of concerts, and the contract to hire it was discharged \citep{taylor1863}. In \emph{Krell v Henry} a room was hired to view the coronation procession of Edward VII, which was postponed when the King fell ill, and the hire contract was held frustrated because the viewing of the procession was the foundation of the contract \citep{krell1903}. In \emph{Davis Contractors v Fareham UDC} the House of Lords emphasised that frustration requires a change so fundamental that performance would be a thing radically different from that undertaken, and that mere hardship or expense is not enough \citep{davis1956}.\index{frustration}

Frustration allocates risks the contract did not address. The Bayesian question is which risks the parties contemplated and priced. An event that both parties regarded as negligible at the time of contracting, such as the postponement of a coronation, was not priced, and the contract's terms say nothing about it. Holding one party to the bargain would impose a loss that neither party anticipated or paid for. The doctrine's narrowness, refusing relief for events that merely made performance more costly and treating foreseeability as weighing against relief, reflects the view that commercial parties may be taken to have priced the risks they could have foreseen. Posner and Rosenfield argued that the efficient rule places the risk of an unforeseen event on the party that could have borne it at least cost, by prevention or insurance \citep{posner1977}. Their analysis supplements the Bayesian one: the prior governs whether the risk was contemplated, and the cost of bearing it governs who should bear it if not.

\section{The Uniform Commercial Code}

American sales law, codified in Article 2 of the Uniform Commercial Code, adopts several rules that are explicit about uncertainty. Section 2-615 of Article 2 excuses a seller from delay or non delivery when performance has been made impracticable by the occurrence of a contingency the non occurrence of which was a basic assumption on which the contract was made, a codification of commercial impracticability, the American counterpart of frustration, for sales. The Code's default rules on warranties, risk of loss and remedies are designed to reflect what most parties would agree to, which is a population prior over contract terms. Karl Llewellyn, the Code's principal architect, wanted rules that reflected commercial practice, and the Code's frequent references to usage of trade and course of dealing direct courts to the commercial context as the source of the prior for interpreting and filling gaps in contracts.

\section{Machines and contracts}

Contract drafting and review are among the tasks to which language models are most often applied. A model can draft clauses, compare a contract with a standard form, identify unusual terms and predict how courts might interpret ambiguous language. In the Bayesian framework, the model supplies estimates of the likelihood $\Prob(T\given m,C)$, how probable it is that parties intending a given meaning would have used this text, and of the prior, how commonly parties in this context agree to a given term. Both are learned from large collections of contracts. The danger, as elsewhere, is that the model's estimates are uncalibrated, or that it smooths over the unusual terms that matter most. The discipline is the same: report probabilities, explain them, and let human judgment decide what the parties agreed.

Smart contracts, programs on blockchains that execute automatically, represent the opposite extreme: contracts whose meaning is their code, with no interpretation. The collapse of the DAO in 2016, when an attacker exploited a feature of its code to drain funds, showed the cost of a contract that does exactly what it says. A legal system that interprets contracts in light of their purpose would not have enforced the exploit; a code that executes literally could not prevent it. The contrast is a reminder that interpretation by inference, with its prior from commercial sense, is a feature of contract law, not a defect.

\section*{Notes and Questions}
\addcontentsline{toc}{section}{Notes and Questions}
\begin{enumerate}[label=\arabic*.]
\item Construct a numerical example of interpretation with two meanings, a likelihood ratio from the text, and a prior from commercial common sense. Show how the result changes as the drafting becomes less precise.
\item In \emph{Arnold v Britton}, the bargain made sense when made but became absurd later. Should a court interpret the clause by the prior at the time of contracting or at the time of the dispute? Why?
\item Explain the officious bystander test as a posterior threshold. What threshold would you choose, and why is it higher than one half?
\item Using the information forcing analysis of \emph{Hadley v Baxendale}, explain why a rule allowing recovery of all losses would raise prices for low loss customers.
\item In \emph{Makdessi}, the clause deprived a seller of deferred consideration if he competed with the business. Assess the clause as an incentive mechanism. What legitimate interest did it protect?
\item Draft a short clause allocating the risk of a pandemic in a supply contract. What probabilities did you assume?
\end{enumerate}

\section*{Further reading}
Lord Hoffmann's speech in \emph{Investors Compensation Scheme} \citep{ics1998}; Ayres and Gertner on default rules \citep{ayres1989}; Posner and Rosenfield on impossibility \citep{posner1977}.
